\subsection{Цель \texttt{compile}}

\subsubsection{Листинг}

\begin{lstlisting}[caption={Часть содержимого файла \texttt{build.xml}},captionpos=b]
<target name="compile">
  <echo message="Compiling the project from ${fse.compile.source}..."/>

  <echo message="Source Java: ${fse.compile.jvm.source}"/>
  <echo message="Target Java: ${fse.compile.jvm.target}"/>

  <delete dir="${fse.compile.target}" />
  <mkdir dir="${fse.compile.target}" />

  <javac srcdir="${fse.compile.source}"
         destdir="${fse.compile.target}"
         source="${fse.compile.jvm.source}"
         target="${fse.compile.jvm.target}"
         classpathref="fse.ivy.path.compile"
         includeantruntime="false">
    <compilerarg value="-proc:none"/>
  </javac>

  <echo message="Done task 'compile'!" />
</target>
\end{lstlisting}

\subsubsection{Комментарии}

Для компиляции исходного кода проекта на Java используется задача \texttt{javac}. Она входит в состав Apache Ant API, поэтому доступна без сторонних библиотек. Компилятору в качестве параметров подаются версии исходного кода и желаемой версии выходных классов. Это нужно, чтобы переиспользовать цель в дальнейшем.

Из других параметров указывается путь к внешним зависимостям проекта \texttt{classpathref} и исключаются Java-зависимости самого Apache Ant, чтобы компиляция была воспроизводимой. 

Перед компиляцией проекта цель пересоздаёт директорию компиляции во избежание ошибок. Это осуществляется задачами \texttt{remove} и \texttt{mkdir}, которые также входят в состав Apache Ant API.

Цель оснащена логгирующими сообщениями, которые реализованы с помощью задачи \texttt{echo} \textit{(доступна в составе Apache Ant API)}.

В выводе цели возникает предупреждение о версии исходного кода Java. Оно гласит, что текущие параметры компиляции не синхронизированы с установленной версией в контейнере. Обычно это означает, что исходный код проекта старше, чем установленная версия JVM. Это никак не повлияет на работоспособность проекта благодаря обратной совместимости.  

\subsubsection{Значения используемых параметров}

\begin{lstlisting}[language=Properties,caption={Часть содержимого файла \texttt{build.properties}},captionpos=b]
fse.compile.source = src
fse.compile.target = build
fse.compile.jvm.source = 17
fse.compile.jvm.target = 17
\end{lstlisting}

\subsubsection{Лог цели}

\begin{lstlisting}[caption={Лог цели \texttt{compile}},language={},captionpos=b]
Buildfile: /project/build.xml

compile:
     [echo] Compiling the project from src...
     [echo] Source Java: 17
     [echo] Target Java: 17
   [delete] Deleting directory /project/build
    [mkdir] Created dir: /project/build
    [javac] Compiling 22 source files to /project/build
    [javac] warning: [options] system modules path not set in conjunction with -source 17
    [javac] 1 warning
     [echo] Done task 'compile'!

BUILD SUCCESSFUL
Total time: 1 second
\end{lstlisting}